\section{Supplementary Methods: implicit mapping, SIJ kinematics, and size scaling}

This Supporting Information documents the computational definitions used to (i) represent subject-specific anatomy via an implicit mapping on a common reference domain and to pull back loads defined on the mapped geometry, (ii) extract SIJ kinematics from deformed geometry, and (iii) compute global size proxies and fit size-scaling models.

\subsection{Implicit geometry mapping and pullback}
For each subject $i$, the FE solve is performed on a common reference domain $\Omega_{ref}$ with coordinates $\mathbf{X}$. Subject-specific geometry is encoded implicitly by a displacement mapping $\phi_i(\mathbf{X})$ such that mapped (subject) coordinates are
\begin{equation}
\mathbf{x}=\boldsymbol{\chi}_i(\mathbf{X})=\mathbf{X}+\phi_i(\mathbf{X}).
\end{equation}
In this repository, subject-specific shape information is provided on the template surface as pointwise displacements relative to the template. The mapping field $\phi_i$ is then constructed by interpolating these displacements to the FE function space (RBF interpolation), yielding an implicit representation of individual geometry on the shared reference mesh. The template (reference) geometry corresponds to $\phi=0$.
Define the deformation gradient and Jacobian of this mapping as
\begin{equation}
\mathbf{F}_i=\mathbf{I}+\nabla \phi_i,
\qquad
J_i=\det(\mathbf{F}_i).
\end{equation}
The linear elastic displacement field $\mathbf{u}_{i\ell}(\mathbf{X})$ for load case $\ell$ is solved in reference coordinates using the standard pullback of gradients by $\mathbf{F}_i^{-1}$ \citep{BonetWood1997}:
\begin{equation}
\boldsymbol{\varepsilon}(\mathbf{w})=\mathrm{sym}\big(\nabla \mathbf{w}\,\mathbf{F}_i^{-1}\big).
\end{equation}

\paragraph{Surface measure and traction pullback.}
Let $\Gamma_{ref}$ denote a reference boundary region with outward unit normal $\mathbf{n}_{ref}$, and let $\Gamma_i=\boldsymbol{\chi}_i(\Gamma_{ref})$ be the corresponding mapped (subject) surface. The surface measure transforms as
\begin{equation}
dS_x = J_{surf}\, dS_X,
\qquad
J_{surf}=J_i\,\big\|\mathbf{F}_i^{-\mathsf{T}}\mathbf{n}_{ref}\big\|.
\end{equation}
Therefore, a physical traction $\mathbf{t}_{phys}$ defined on $\Gamma_i$ contributes
\begin{equation}
\int_{\Gamma_i} \mathbf{t}_{phys}\cdot \mathbf{v}\, dS_x
=
\int_{\Gamma_{ref}} \mathbf{t}_{phys}\cdot \mathbf{v}\; J_{surf}\, dS_X.
\end{equation}

\paragraph{Total-force standardization.}
To impose a prescribed resultant force vector $\mathbf{F}^{tot}$ on the mapped surface, a constant traction is chosen as
\begin{equation}
\mathbf{t}_{phys}=\frac{\mathbf{F}^{tot}}{A_{phys}},
\qquad
A_{phys}=\int_{\Gamma_{ref}} J_{surf}\, dS_X,
\end{equation}
which ensures $\int_{\Gamma_i} \mathbf{t}_{phys}\, dS_x=\mathbf{F}^{tot}$. (Pressure loads use the corresponding Piola transform $J_i \mathbf{F}_i^{-\mathsf{T}}\mathbf{n}_{ref}$.)

\paragraph{Deformed coordinates for point clouds.}
For a reference point $\mathbf{X}$ on the template surface, the deformed coordinate used for kinematics is
\begin{equation}
\mathbf{x}_{i\ell}(\mathbf{X})=\mathbf{X}+\phi_i(\mathbf{X})+\mathbf{u}_{i\ell}(\mathbf{X}).
\end{equation}
In practice, $\phi_i$ and $\mathbf{u}_{i\ell}$ are evaluated at surface points by interpolation from the FE degrees of freedom.
In the present workflow, this evaluation uses radial basis function (RBF) interpolation constructed on FE degree-of-freedom coordinates.

\subsection{Notation and inputs}
For each subject $i$ and load case $\ell$, the finite-element solution defines a deformed pelvis geometry obtained by evaluating $\mathbf{x}_{i\ell}(\mathbf{X})$ on the template surface points (one-to-one correspondence). In this repository, the reference surface mesh (``ref'') is the common template in its reference configuration, and the deformed surface mesh (``def'') is obtained by applying the subject mapping $\phi_i$ and the load-specific displacement $\mathbf{u}_{i\ell}$ to the same point set.
After splitting the surface into bodies (left ilium, right ilium, sacrum), denote the reference point clouds as
$\{\mathbf{X}^{LI}_j\}$, $\{\mathbf{X}^{RI}_j\}$, and $\{\mathbf{X}^{S}_j\}$, and the corresponding deformed point clouds as
$\{\mathbf{x}^{LI}_j\}$, $\{\mathbf{x}^{RI}_j\}$, and $\{\mathbf{x}^{S}_j\}$.

All vectors are 3D column vectors.
Rotation matrices are $3\times 3$ and are orthogonal with determinant $+1$.
Translations are 3D vectors.
Angles are reported in degrees and translations in mm.

\subsection{Pelvis-aligned coordinate frame (ML/AP/CC)}
All SIJ outputs are expressed in a pelvis-aligned orthonormal frame whose columns are the unit axes
$\hat{\mathbf{x}}$ (ML), $\hat{\mathbf{y}}$ (AP), and $\hat{\mathbf{z}}$ (CC).
The frame is computed from the reference (undeformed) pelvis point cloud.

\paragraph{Origin.}
Let $\mathbf{o}$ be the centroid of all reference surface points:
\begin{equation}
\mathbf{o}=\frac{1}{N}\sum_{j=1}^{N}\mathbf{X}_j.
\end{equation}

\paragraph{ML axis.}
Compute the sample covariance matrix of the centered point cloud and its eigen-decomposition.
Let $\hat{\mathbf{x}}$ be the eigenvector associated with the largest eigenvalue (the dominant PCA axis), oriented so that its global $x$ component is nonnegative (a deterministic sign convention):
\begin{equation}
\hat{\mathbf{x}}=\arg\max_{\|\mathbf{v}\|=1}\mathbf{v}^\mathsf{T}\,\mathrm{cov}(\mathbf{X}_j-\mathbf{o})\,\mathbf{v}.
\end{equation}

\paragraph{CC and AP axes.}
Project the centered point cloud into the plane orthogonal to $\hat{\mathbf{x}}$:
\begin{equation}
\tilde{\mathbf{X}}_j=(\mathbf{X}_j-\mathbf{o})-\big((\mathbf{X}_j-\mathbf{o})\cdot\hat{\mathbf{x}}\big)\hat{\mathbf{x}}.
\end{equation}
The dominant direction in this plane is extracted by PCA/SVD and then orthogonalized against $\hat{\mathbf{x}}$ to obtain $\hat{\mathbf{z}}$.
The remaining axis is defined to ensure a right-handed frame:
\begin{equation}
\hat{\mathbf{y}}=\frac{\hat{\mathbf{z}}\times\hat{\mathbf{x}}}{\|\hat{\mathbf{z}}\times\hat{\mathbf{x}}\|},
\qquad
\hat{\mathbf{z}}=\frac{\hat{\mathbf{x}}\times\hat{\mathbf{y}}}{\|\hat{\mathbf{x}}\times\hat{\mathbf{y}}\|}.
\end{equation}

\paragraph{Left--right sign convention.}
To enforce that positive ML points from left to right, the ML axis is flipped (with the AP axis simultaneously flipped to preserve right-handedness) if the left ilium mean projects more positively than the right ilium mean along $\hat{\mathbf{x}}$:
\begin{equation}
(\bar{\mathbf{X}}^{LI}-\mathbf{o})\cdot\hat{\mathbf{x}} > (\bar{\mathbf{X}}^{RI}-\mathbf{o})\cdot\hat{\mathbf{x}}
\;\Rightarrow\;
\hat{\mathbf{x}}\leftarrow-\hat{\mathbf{x}},\;
\hat{\mathbf{y}}\leftarrow-\hat{\mathbf{y}}.
\end{equation}

Collect axes into the orthonormal frame matrix
\begin{equation}
\mathbf{P}=\begin{bmatrix}\hat{\mathbf{x}} & \hat{\mathbf{y}} & \hat{\mathbf{z}}\end{bmatrix},
\end{equation}
which maps world vectors into pelvis coordinates by $\mathbf{v}_{pel}=\mathbf{P}^\mathsf{T}\mathbf{v}$.

\subsection{Removing global affine motion using the sacrum}
To focus on SIJ-relative motion, any global affine deformation of the entire pelvis is removed using a sacrum-based fit.
Let $\mathbf{X}^{S}_j$ and $\mathbf{x}^{S}_j$ be corresponding sacrum points in reference and deformed meshes.
We estimate an affine map $(\mathbf{A},\mathbf{b})$ by least squares:
\begin{equation}
(\mathbf{A},\mathbf{b})
=
\arg\min_{\mathbf{A},\mathbf{b}}
\sum_j \big\| \mathbf{A}\mathbf{X}^{S}_j+\mathbf{b}-\mathbf{x}^{S}_j \big\|^2.
\end{equation}
The inverse affine map is then applied to all deformed points (sacrum and both ilia) to obtain affine-corrected deformed coordinates:
\begin{equation}
\mathbf{x}'=\mathbf{A}^{+}\,(\mathbf{x}-\mathbf{b}),
\end{equation}
where $\mathbf{A}^{+}$ denotes the Moore--Penrose pseudoinverse (robust to mild ill-conditioning).

\subsection{Defining a side-specific SIJ region of interest (ROI)}
SIJ rigid fits are computed over surface ROIs intended to approximate the articular contact regions.
On each side $s\in\{L,R\}$ we first select candidate points by the sign of the ML coordinate in the reference frame:
\begin{equation}
s=L:\ (\mathbf{X}-\mathbf{o})\cdot\hat{\mathbf{x}}<0,
\qquad
s=R:\ (\mathbf{X}-\mathbf{o})\cdot\hat{\mathbf{x}}>0.
\end{equation}
Within each side, ROIs are defined by proximity between the sacrum and ilium surfaces in the reference configuration.
For a sacrum point $\mathbf{p}$, define its distance to the opposite bone surface as
\begin{equation}
d(\mathbf{p},\mathcal{B})=\min_{\mathbf{q}\in\mathcal{B}}\|\mathbf{p}-\mathbf{q}\|.
\end{equation}
Given a gap tolerance $g=\SI{5}{mm}$, the sacral ROI on side $s$ is
\begin{equation}
\mathcal{S}_s=\{\mathbf{X}^{S}_j:\ d(\mathbf{X}^{S}_j,\{\mathbf{X}^{I_s}_k\})\le g\},
\end{equation}
and the corresponding iliac ROI is
\begin{equation}
\mathcal{I}_s=\{\mathbf{X}^{I_s}_k:\ d(\mathbf{X}^{I_s}_k,\{\mathbf{X}^{S}_j\})\le g\}.
\end{equation}
If either ROI contains fewer than $N_{min}=80$ points, the $N_{min}$ closest points by $d(\cdot)$ are used instead (a deterministic fallback to avoid empty/small ROIs).

\subsection{Rigid-body fit on each bone ROI (trimmed Kabsch)}
For each side $s$ and each bone (sacrum and ilium), a rigid transform $(\mathbf{R},\mathbf{t})$ is fit that maps reference ROI points to affine-corrected deformed ROI points.
Given corresponding point sets $\{\mathbf{p}_k\}$ and $\{\mathbf{q}_k\}$, we solve
\begin{equation}
(\mathbf{R},\mathbf{t})
=
\arg\min_{\mathbf{R}\in SO(3),\,\mathbf{t}}
\sum_{k}\|\mathbf{R}\mathbf{p}_k+\mathbf{t}-\mathbf{q}_k\|^2.
\end{equation}
The Kabsch solution is used.
Let $\bar{\mathbf{p}}$ and $\bar{\mathbf{q}}$ be the point-set centroids and define centered coordinates
$\mathbf{p}^c_k=\mathbf{p}_k-\bar{\mathbf{p}}$, $\mathbf{q}^c_k=\mathbf{q}_k-\bar{\mathbf{q}}$.
The cross-covariance is
\begin{equation}
\mathbf{C}=\sum_k \mathbf{p}^c_k(\mathbf{q}^c_k)^\mathsf{T}.
\end{equation}
With SVD $\mathbf{C}=\mathbf{U}\mathbf{\Sigma}\mathbf{V}^\mathsf{T}$, the rotation is
\begin{equation}
\mathbf{R}=\mathbf{V}\mathbf{D}\mathbf{U}^\mathsf{T},
\qquad
\mathbf{D}=\mathrm{diag}\big(1,1,\det(\mathbf{V}\mathbf{U}^\mathsf{T})\big),
\end{equation}
and the translation is
\begin{equation}
\mathbf{t}=\bar{\mathbf{q}}-\mathbf{R}\bar{\mathbf{p}}.
\end{equation}

To reduce sensitivity to outliers, the fit is \emph{trimmed}:
residuals $e_k=\|\mathbf{R}\mathbf{p}_k+\mathbf{t}-\mathbf{q}_k\|$ are computed,
the smallest $(1-\tau)$ fraction is retained ($\tau=0.1$), and the rigid fit is refit on the retained subset.
This trimming--refit procedure is run for two iterations.

\subsection{Relative SIJ rotation and translation}
Let $(\mathbf{R}^{S}_s,\mathbf{t}^{S}_s)$ and $(\mathbf{R}^{I}_s,\mathbf{t}^{I}_s)$ be the rigid transforms for the sacrum and the ilium on side $s$.
The ilium motion relative to the sacrum is expressed by
\begin{equation}
\mathbf{R}^{rel}_s=(\mathbf{R}^{S}_s)^\mathsf{T}\mathbf{R}^{I}_s,
\qquad
\mathbf{t}^{rel}_s=(\mathbf{R}^{S}_s)^\mathsf{T}\big(\mathbf{t}^{I}_s-\mathbf{t}^{S}_s\big).
\end{equation}

\paragraph{Euler angles (NUT/AP/CC).}
To report rotations about the pelvis axes, we express the relative rotation in the pelvis frame:
\begin{equation}
\mathbf{R}^{pel}_s=\mathbf{P}^\mathsf{T}\mathbf{R}^{rel}_s\mathbf{P}.
\end{equation}
Euler XYZ angles are then extracted as
\begin{equation}
\alpha_x=\mathrm{atan2}\!\big(\mathbf{R}^{pel}_{32},\mathbf{R}^{pel}_{33}\big),
\quad
\beta_y=\arcsin\!\big(-\mathbf{R}^{pel}_{31}\big),
\quad
\gamma_z=\mathrm{atan2}\!\big(\mathbf{R}^{pel}_{21},\mathbf{R}^{pel}_{11}\big),
\end{equation}
with standard singularity handling when $|\mathbf{R}^{pel}_{31}|\approx 1$ (gimbal-lock case).
Angles are converted to degrees and reported as NUT ($\alpha_x$), AP-axis rotation ($\beta_y$), and CC-axis rotation ($\gamma_z$).

\paragraph{Translation (ML/AP/CC).}
We define a representative SIJ contact point as the centroid of the sacral ROI in the \emph{reference} configuration,
$\mathbf{p}_s=\mathrm{mean}(\mathcal{S}_s)$.
The relative displacement between ilium and sacrum at this point (in world coordinates) is computed as
\begin{equation}
\mathbf{d}^{world}_s=
\big(\mathbf{R}^{I}_s-\mathbf{R}^{S}_s\big)\mathbf{p}_s+\big(\mathbf{t}^{I}_s-\mathbf{t}^{S}_s\big),
\end{equation}
and expressed in pelvis coordinates by
\begin{equation}
\mathbf{d}^{pel}_s=\mathbf{P}^\mathsf{T}\mathbf{d}^{world}_s=
\begin{bmatrix}d_{ML}\\ d_{AP}\\ d_{CC}\end{bmatrix}_s.
\end{equation}
These components are reported (mm) as $dx_{ML}$, $dy_{AP}$, and $dz_{CC}$ for the left and right SIJ.

\subsection{Derived subject-level SIJ metrics used in the manuscript}
For each subject $i$, load $\ell$, and side $s\in\{L,R\}$, define the rotation and translation component vectors
\begin{equation}
\boldsymbol{\theta}_{i\ell s}=
\begin{bmatrix}\alpha_x\\ \beta_y\\ \gamma_z\end{bmatrix}_{i\ell s},
\qquad
\mathbf{d}_{i\ell s}=
\begin{bmatrix}d_{ML}\\ d_{AP}\\ d_{CC}\end{bmatrix}_{i\ell s}.
\end{equation}
The side-specific magnitudes are the Euclidean norms
\begin{equation}
r_{i\ell s}=\|\boldsymbol{\theta}_{i\ell s}\|_2,
\qquad
t_{i\ell s}=\|\mathbf{d}_{i\ell s}\|_2.
\end{equation}
Primary outcomes are the bilateral means of these magnitudes:
\begin{equation}
\mathrm{rot\_mag\_deg}_{i\ell}=\tfrac{1}{2}(r_{i\ell L}+r_{i\ell R}),
\qquad
\mathrm{trans\_mag\_mm}_{i\ell}=\tfrac{1}{2}(t_{i\ell L}+t_{i\ell R}).
\end{equation}
Directional component summaries are bilateral means of absolute component magnitudes, e.g.
\begin{equation}
\mathrm{ap\_trans\_abs\_mm}_{i\ell}=\tfrac{1}{2}\big(|d_{AP,i\ell L}|+|d_{AP,i\ell R}|\big),
\end{equation}
with analogous definitions for ML and CC translation and for the rotational components (nutation/AP/CC).
Left--right asymmetry magnitudes are defined on the scalar magnitudes:
\begin{equation}
\mathrm{lr\_rot\_asym\_deg}_{i\ell}=|r_{i\ell L}-r_{i\ell R}|,
\qquad
\mathrm{lr\_trans\_asym\_mm}_{i\ell}=|t_{i\ell L}-t_{i\ell R}|.
\end{equation}

\subsection{Global size proxies and size-scaling model}
\paragraph{Mapped volume, surface area, and scale.}
Global size proxies were computed from the mapped pelvis geometry implied by $\phi_i$. Using the implicit mapping $\boldsymbol{\chi}_i$, the mapped total pelvic volume $V_i$ (cm$^3$) and total pelvic surface area $S_i$ (cm$^2$) can be expressed as
\begin{equation}
V_i=\int_{\Omega_{ref}} J_i(\mathbf{X})\,d\mathbf{X},
\qquad
S_i=\int_{\Gamma_{ref}} J_{surf}(\mathbf{X})\,dS_X,
\end{equation}
with $J_i=\det(\mathbf{F}_i)$ and $J_{surf}$ defined above.
An isotropic global scale factor is defined by
\begin{equation}
\mathrm{scale}_i=\left(\frac{V_i}{V_{template}}\right)^{1/3},
\end{equation}
where $V_{template}$ is the total volume of the template pelvis.

\paragraph{Size-scaling model.}
For each load case $\ell$ and each primary SIJ outcome $Y_{i\ell}$ (rotation magnitude and translation magnitude), we fit a log--log model
\begin{align}
\log(Y_{i\ell}) &=
\beta_0
+\beta_1\log(V_i)
+\beta_2\,\mathrm{sex}_i
+\beta_3\,\mathrm{age}_{z,i}
+\beta_4\,\mathrm{AP}_{z,i}
+\beta_5\,\mathrm{Biischiadic}_{z,i}
+\beta_6\,\mathrm{Subpubic}_{z,i} \nonumber\\
&\quad
+\beta_7\log(V_i)\,\mathrm{sex}_i
+\varepsilon_{i\ell},
\end{align}
where $\mathrm{sex}_i=1$ for female and $0$ for male, and $x_{z}$ denotes $z$-scoring
$x_z=(x-\bar{x})/s_x$ using the sample standard deviation ($ddof=1$).
Models are estimated by OLS with HC3 heteroskedasticity-consistent covariance.

\paragraph{Exponent extraction and confidence intervals.}
In the above parameterization, the male scaling exponent is $b_M=\beta_1$ and the female scaling exponent is $b_F=\beta_1+\beta_7$.
Wald-type 95\% CIs use the model covariance matrix:
\begin{equation}
b_M \pm 1.96\sqrt{\mathrm{Var}(\beta_1)},
\qquad
b_F \pm 1.96\sqrt{\mathrm{Var}(\beta_1)+\mathrm{Var}(\beta_7)+2\,\mathrm{Cov}(\beta_1,\beta_7)}.
\end{equation}
The sex-by-size slope interaction is tested by the $p$ value of $\beta_7$, with Benjamini--Hochberg false-discovery control applied across the interaction tests.
